\subsection{合成列的利用}

设 A 是一个环。设 E 是一个有限长度的 A-模，S 是一个单 A-模。由若尔当–赫尔德定理（I, §4, No.~7, 第 43 页，定理 6），E 的若尔当–赫尔德列中同构于 S 的商的个数不依赖于该序列。我们将它记作 \(\ell_{\mathrm{S}}(\mathrm{E})\)，并称之为 S 在 E 中的重数。A-模 E 的支集是使 \(\ell_{\mathrm{S}}(\mathrm{E}) \neq 0\) 的单 A-模 S 的类所成之集。当 E 是半单的且具有有限长度时，整数 \(\ell_{\mathrm{S}}(\mathrm{E})\) 是 E 的 S 型同型分量的长度 {[}E : S{]}（VIII, 第 72 页），且支集的概念与 VIII, 第 66 页中引入的概念一致。

引理 1. —— 设 E、E' 与 E'\,' 是有限长度的 A-模，且

\[
0 \longrightarrow \mathrm{E} ^ {\prime} \stackrel {i} {\longrightarrow} \mathrm{E} \stackrel {p} {\longrightarrow} \mathrm{E} ^ {\prime \prime} \longrightarrow 0
\]

是一个正合列。我们有 \(\ell_{\mathrm{S}}(\mathrm{E}) = \ell_{\mathrm{S}}(\mathrm{E}') + \ell_{\mathrm{S}}(\mathrm{E}'')\)。

设 \(\Sigma'\) 与 \(\Sigma''\) 分别是 \(i(\mathrm{E}')\) 与 \(\mathrm{E}/i(\mathrm{E}')\) 的若尔当–赫尔德列。存在一个若尔当–赫尔德列 \(\Sigma\)（\(\mathrm{E}\) 的），其商序列可由 \(\Sigma\) 与 \(\Sigma'\) 的商序列并置得到（I, §4, No.~8, 第 44 页）。

命题 7. —— 设 C 是模的类的一个遗传集，使得每个 C 型模都具有有限长度。设 S 是属于 C 的单模的类所成之集。则族 \(([S]_{\mathcal{C}})_{S \in \mathcal{S}}\) 是 Z-模 \(K(\mathcal{C})\) 的一个基，且我们有

\[
[ \mathrm{E} ] _ {\mathcal {C}} = \sum_ {\mathrm{S} \in \mathcal {S}} \ell_ {\mathrm{S}} (\mathrm{E}) [ \mathrm{S} ] _ {\mathcal {C}}\tag{8}
\]

（对每个 \(\mathcal{C}\) 型模 E）。

公式 (8) 由将命题 3 应用于 E 的一个若尔当–赫尔德列得出。由引理 1，对 S 的每个元素 S，存在一个 Z-线性映射 \(\varphi_{\mathrm{S}}\)（从 \(\mathrm{K}(\mathcal{C})\) 到 Z），使得对每个 C 型模 E 有 \(\varphi_{\mathrm{S}}([\mathrm{E}]_{\mathcal{C})} = \ell_{\mathrm{S}}(\mathrm{E})\)。特别地，我们有 \(\varphi_{\mathrm{S}}([\mathrm{S}]_{\mathcal{C}}) = 1\)，且 \(\varphi_{\mathrm{S}}([\mathrm{S}']_{\mathcal{C}}) = 0\) 对 S 中每个 \(S' \neq S\) 成立。由此得出，形如 \([S]_{\mathcal{C}}\)（\(S \in S\)）的元素在 Z 上线性无关；由公式 (8)，这些元素生成 \(\mathrm{K}(\mathcal{C})\)。

推论. —— 设 E 与 F 是 C 型半单模。模 E 同构于 F 当且仅当我们有 \([E]_{C} = [F]_{C}\)（在 \(K(\mathcal{C})\) 中）。

事实上，我们有 \([E]_{\mathcal{C}} = \sum_{S \in \mathcal{S}} \ell_{S}(E)[S]_{\mathcal{C}}\)，对 F 也有类似的公式，且 E 同构于 F 当且仅当我们有 \(\ell_{S}(E) = \ell_{S}(F)\)（对每个 \(S \in \mathcal{S}\)）（VIII, 第 72 页）。

注. —— 集 \(\mathrm{K}(\mathcal{C})^{+}\) 是 \(\mathrm{K}(\mathcal{C})\) 中由族 \(([S]_{\mathcal{C}})_{S\in\mathcal{S}}\) 生成的子幺半群。

